\chapter{稳态：身体不是保持不变，而是在不停纠偏}

\begin{kaichang}
找一支体温计，夹在腋下五分钟。不管你是在盛夏的午后测，还是在寒冬的清晨测；不管你刚喝完热汤，还是刚从冷气房里走出来——读数大概率都落在 $36\,^{\circ}\mathrm{C}$ 到 $37\,^{\circ}\mathrm{C}$ 之间那窄窄的一格里。

这件事细想起来有点奇怪。外面的气温一年能差出 $40\,^{\circ}\mathrm{C}$，你吃的东西忽冷忽热，你一会儿冲刺跑一百米、一会儿躺平睡觉，身体内部的"锅炉"功率时大时小。可体温计上的数字几乎不动。是你的身体真的一动不动地"保持"着体温吗？

先猜一个答案。这一章结束时你会发现：真相恰好相反——你的身体内部其实每分每秒都在"出错"，体温、血糖、酸碱度、水分全在波动，体温计上的稳定，是靠一整支不停加班的"纠偏队伍"硬拉回来的。稳定不是不动，而是动了立刻被掰回来。
\end{kaichang}

\section{那些"不听话"的数字}

先只看能测量的东西。用体温计、血糖仪、血压计和一张化验单，你能从自己身上读到一串数字：

\begin{itemize}[itemsep=2pt]
\item 体温：大约 $36.1\,^{\circ}\mathrm{C}$ 到 $37.2\,^{\circ}\mathrm{C}$，一天之内有不到 $1\,^{\circ}\mathrm{C}$ 的起伏，清晨低、傍晚高。
\item 空腹血糖：每升血液里约 $0.7$ 到 $1.0$ 克葡萄糖。饭后会升，两小时内又落回这个区间。
\item 血液的酸碱度（pH）：$7.35$ 到 $7.45$。这个范围窄得吓人——偏出 $0.4$ 就有生命危险。
\item 血钠浓度：每升约 $135$ 到 $145$ 毫摩尔，口渴和多尿会围着它打转。
\item 血氧饱和度：健康人在平原地区通常高于 $95\,\%$。
\end{itemize}

注意两件事。第一，这些数字都有\textbf{范围}，不是钉死的点。第二，你每天都在往这些范围外面"捣乱"：一顿甜饮料把血糖往上顶，一次八百米跑让肌肉狂产乳酸和热量，一顿火锅灌进去大把的盐。可几个小时后，数字又回来了。

这有点像走钢丝的人。远远看去，他稳稳地停在钢丝上；凑近看，他的身体其实一直在左右晃，手里的平衡杆没有一刻不在调整。你的身体就是那位走钢丝的人，只不过它同时走在五六根钢丝上，而且一走就是一辈子。

\section{把镜头推进去：细胞泡在什么里面}

为什么要这么紧张地纠偏？答案得回到粒子层去。

第 57 章说过，你的几十万亿个细胞不是挤成一团的实心砖，而是泡在一层薄薄的液体里——血液、组织液，合起来叫\textbf{内环境}。每个细胞的"吃穿用度"全靠这汪液体送货上门：葡萄糖、氧气、各种离子从它里面取，二氧化碳和代谢废物往它里面扔。

而细胞内部是一整套精密的化学工厂。第 41 章到第 43 章讲过的酶，个个是挑剔的蛋白质：温度高个几度，形状就开始松动，活性直线下降；再高，直接变性，像煮熟的蛋清一样再也回不去（第 44 章）。pH 偏一点，蛋白质表面带的电荷就变了，该结合的底物结合不上。钠离子、钾离子、钙离子的浓度更是神经和肌肉的命根子——第 58 章讲信号传递时见过，神经冲动本质上就是钠钾离子跨膜的浪潮，浓度不对，心跳都会乱套。

换句话说：\textbf{细胞里的化学反应网络，是在一份非常苛刻的"环境规格书"下调试出来的}。内环境偏离规格，工厂就降速、出错、停产。发烧到 $41\,^{\circ}\mathrm{C}$ 会昏迷，不是因为身体"太娇嫩"，而是因为上亿个酶分子同时开始掉链子。

所以那些数字必须被盯住。可谁来盯？盯住了又怎么往回拉？

\section{纠偏的通用图纸：负反馈}

大自然（其实是演化，第 83 章以后的故事）给这个问题提供了一套反复使用的图纸，叫\textbf{负反馈}。它只有三个角色，你家里就有一台照着这套图纸工作的机器——空调：

\begin{itemize}[itemsep=2pt]
\item \textbf{传感器}：空调的测温探头；身体里的温度感受器、化学感受器，随时报告"现在是多少"。
\item \textbf{控制中心}：空调的电路板；身体里的下丘脑（脑里一块指甲盖大小的区域）、胰岛细胞等，把报告值和\textbf{设定值}比较，算出偏差。
\item \textbf{执行器}：空调的压缩机；身体里的汗腺、肌肉、肝、肾、血管，负责把数值往回拉。
\end{itemize}

关键在那个"负"字：偏高就往下压，偏低就往上抬，纠偏的方向永远和偏差相反。拿体温走一遍完整流程——夏天你站在太阳底下，体温往上爬了 $0.5\,^{\circ}\mathrm{C}$：皮肤和下丘脑的温度感受器察觉升温，下丘脑发出指令，汗腺开闸出汗、皮肤血管舒张，把热量往体表送、靠蒸发带走；体温回落，指令撤销。冬天反过来：皮肤血管收缩减少散热，肌肉不自主地打寒战——寒战不是"冻得发抖"这么简单，它是骨骼肌在空转做功，把 ATP 的能量直接变成热量，一台临时启动的"人肉暖气"。

血糖走同一套图纸，只是控制中心和执行器换了角色。饭后血糖升高，胰岛 $\beta$ 细胞自己就是"传感器加控制中心"的合体，立刻释放胰岛素，催促肝、肌肉、脂肪细胞把葡萄糖收进仓库（第 55 章讲过这条代谢指令的细节）；饿的时候血糖下滑，胰岛 $\alpha$ 细胞放出胰高血糖素，命令肝把糖原拆回葡萄糖放回血液。两套人马方向相反、轮流上岗，血糖就被夹在中间那条窄带里。

\begin{qiaoliang}
这套纠偏系统勤快到什么程度？算一笔账。你的全身血液大约 $5$ 升，按每升 $1$ 克葡萄糖的上限算，\textbf{你全身血管里漂着的葡萄糖总共只有约 $5$ 克}——还不到一小块方糖。而一听可乐含糖约 $35$ 克，是你全部血糖的 $7$ 倍。也就是说，喝完一听可乐，如果没有纠偏，你的血糖浓度理论上会被顶到正常值的 $8$ 倍。但饭后两小时再测，数字基本回来了：几十克葡萄糖在几十分钟内被分送进肝和肌肉的仓库。再看热量：你安静坐着时，全身的细胞呼吸（第 53 章）功率约 $80$ 到 $100$ 瓦，和一盏老式白炽灯泡相当。假设热量完全散不出去，身体按约 $245$ 千焦每摄氏度的"热容量"升温（体重 $70$ 千克，比热接近水），一小时就能升温约 $1.4\,^{\circ}\mathrm{C}$——半天之内就会热到危险线。你之所以没有"发烧发到熟"，全靠散热系统时时刻刻在跟产热打平手。
\end{qiaoliang}

负反馈追求的不是"停在某一点"，而是"在设定值附近小幅震荡"。真实的血糖曲线、体温曲线都是波浪线，只是波幅被死死摁住。标题里那句话现在可以说完整了：\textbf{稳态不是保持不变，而是偏差一冒头就被纠正}。

\section{也有"火上浇油"的时候：正反馈}

负反馈是纠偏，那有没有反着来的——偏差越大、越要加码？有，叫\textbf{正反馈}，身体把它用在"必须一鼓作气干到底"的场合。

最典型的是分娩：子宫收缩压迫胎儿，胎儿压迫宫颈，宫颈的神经信号传回大脑，大脑下令分泌更多催产素，催产素让子宫收缩得更厉害——一轮推着一轮，直到孩子出生，这个"越缩越缩"的循环才被物理性地终止。血液凝固也类似：凝血因子一旦被激活，会激活更多同伴，像多米诺骨牌一样迅速放大，赶在失血过多之前把口子堵上。

注意正反馈的两个特征：它是\textbf{局部的、有终点的}。分娩以胎儿娩出收场，凝血以伤口愈合刹车。如果让它像负反馈那样长期当家，结果就是把一个参数推到天边——那恰恰是失控。所以身体的默认模式是负反馈，正反馈只是被严格圈养的"突击队"。

\section{快马与慢信使：神经和内分泌怎么分工}

纠偏指令靠什么送达？两条线路，性格完全不同，恰好互补（第 58 章细讲过信号分子和受体，这里只看分工）。

\textbf{神经}是点对点专线：电脉冲沿着神经纤维以每秒几十米的速度冲刺，毫秒级到达，指哪打哪，收回也快。寒战、血管收缩、心跳加速，都走这条线——需要立刻响应的，全是它。

\textbf{内分泌}是广播慢递：激素从腺体释放进血液，随血流漂遍全身，几秒到几分钟才送达，而且谁都能"收到"，但只有长着对应受体的细胞才会"拆信"（第 58 章的锁钥模型）。胰岛素、胰高血糖素、抗利尿激素、甲状腺激素，走的都是这条线。慢是慢，但胜在覆盖广、续航久——血糖、水盐这种需要全身几十种细胞协同调整的事，专线根本拉不过来。

一个反应要快，一个局面要稳，两套系统在下丘脑那里会师：下丘脑一边接收全身的温度、渗透压报告，一边既能发神经指令，又能指挥脑垂体这个"激素总调度室"。它就是你身体里那间 $24$ 小时不下班的控制室。

\section{化学家的地盘：肾、肺、肝}

体温、血糖说得多了，其实纠偏最"化学"的部分藏在三个器官里，它们管的是内环境的\textbf{成分}。

\textbf{肺}管两种气体的进出：二氧化碳排多了血液偏碱，排少了偏酸。血液的酸碱缓冲体系（下一节细讲）顶在前面，肺则在后面做粗调——你跑步时呼吸不自觉加快加深，不是"喘不上气"这么简单，是在加速倒掉细胞呼吸猛产出来的 \chem{CO_2}。脖子上的颈动脉里还藏着专门监测血氧和 \chem{CO_2} 浓度的化学感受器，一旦血氧下滑（比如上高原、剧烈运动），它们立刻通过神经催促呼吸肌加班，把氧气摄入拉回来。

\textbf{肝}是化学品集散中心兼解毒车间：葡萄糖在这里存取（第 55 章），多余的氨基酸在这里拆解，拆下来的有毒氨在这里被转成毒性低得多的尿素，交给肾脏带走；酒精、药物也大多在这里被改造失活。

\textbf{肾}是最不知疲倦的质检员：每天把全身血液过滤几十遍——约 $180$ 升原尿，比两大桶桶装水还多——然后把 $99\,\%$ 的水、全部葡萄糖、够用的盐回收回血液，只把多余的水、盐、尿素等"尾货"浓缩成一两升尿液排出。喝了一大杯水，尿很快变多变稀；吃了一顿咸的，尿变浓、你还会口渴——这都是肾和渴觉在联手把血钠浓度摁回那条 $135$ 到 $145$ 的窄带里。抗利尿激素就是这条回路的慢信使：血液变浓时下丘脑释放它，命令肾脏"多回收水"。

三个器官的分工可以一句话记住：肺管气体，肝管转化，肾管进出。内环境这汪水，就是它们仨 $24$ 小时联合维护的。

\section{符号时间：血液怎样"吸收"酸碱冲击}

现在可以给 pH 稳态写出符号了。血液里最重要的缓冲系统是一对搭档：二氧化碳和碳酸氢根。它们之间存在一个可逆平衡：

\[\chem{CO_2} + \chem{H_2O} \rightleftharpoons \chem{H_2CO_3} \rightleftharpoons \chem{H^{+}} + \chem{HCO_3^{-}}\]

回忆第 28 章的平衡移动原理：这个式子像一根化学弹簧。血液里 \chem{H^{+}} 突然增多（比如运动产生的乳酸解离出氢离子），平衡向左移，多余的 \chem{H^{+}} 被 \chem{HCO_3^{-}} 接住、变成 \chem{CO_2}，再被肺呼出去；\chem{H^{+}} 变少时平衡向右移，补上一批。只要这对搭档没耗尽，血液的 pH 就被弹簧拉在 $7.35$ 到 $7.45$ 之间。

注意式子两头的"外援"：\chem{CO_2} 那端连着肺（呼多呼少直接拉动平衡），\chem{HCO_3^{-}} 那端连着肾（肾能回收或排出碳酸氢根，只是动作慢，以小时计）。化学平衡顶秒级的第一波冲击，肺做分钟级的调整，肾做小时级的收尾——同一个变量，三道时间尺度不同的防线。这是稳态系统的通用设计：\textbf{快的先扛，慢的善后}。

\section{科学家怎样知道：从"内环境"到"稳态"}

\begin{zhengju}
今天听起来理所当然的"身体维持内部环境稳定"，在 19 世纪中叶还是个没人想到的问题。当时的默认想法是：血液里的东西直接来自吃喝——血糖高是因为刚吃了糖，如此而已。

法国生理学家克洛德·贝尔纳偏要较这个真。他把狗分成两组，一组照常喂食，一组饿着，然后比较两组动物肝脏和血液里的糖。按照"糖全来自食物"的想法，饿狗的肝脏里不该有糖。结果相反：\textbf{饿着的时候，肝脏照样往血液里放糖}。进一步的追踪让他确认，肝能把一种储存物质（后来命名的糖原）转化成葡萄糖释放进血。这直接推翻了"血糖只是食物的倒影"——身体自己在\textbf{生产}血糖，生产量还随需要而变。

贝尔纳由此提出一个更大的想法：多细胞生物的真正生存环境不是体外那个风霜雨雪的世界，而是体内那汪液体（他称之为"内环境"）。他在 1865 年的讲义里写下那句名言：内环境的恒定，是生命自由独立的前提。注意他凭什么敢这么说——不是哲学感概，而是一连串"饿狗的肝里为什么还有糖"式的实测。

半个世纪后，美国生理学家坎农接手了这个问题：恒定是怎么实现的？他和学生做了大量实验，比如刺激动物的交感神经，观测血糖、心跳、肾上腺分泌的连锁变化，发现身体对"偏差"的响应总是朝着把它压回去的方向。1926 年，他给这套机制正式起名 homeostasis（稳态）——希腊语里"相似"加"停留"，特意没用"完全相同"：他想强调的恰恰是"在一个范围内被拉住"，而不是纹丝不动。再后来，控制论兴起，工程师们发现恒温器、巡航定速和坎农的稳态是同一张图纸——负反馈。生理学和控制工程在两个毫不相干的领域画出了同一个圈，这类巧合通常意味着：抓到的是真规律。
\end{zhengju}

\section{发烧：不是锅炉失控，是有人调高了设定值}

用稳态的眼光重新看发烧，结论会出乎你的意料。

发烧时体温升到 $39\,^{\circ}\mathrm{C}$，这\textbf{不是}散热系统失灵——恰恰相反，散热和产热系统都在正常工作，只是它们维护的"设定值"被人为调高了。证据就在你的主观感觉里：发烧刚开始体温往上爬的那阵子，你会觉得\textbf{冷}，裹被子还打寒战。冷？体温明明在升！这正说明下丘脑的设定值已经先一步调到 $39\,^{\circ}\mathrm{C}$，而你的实际体温还"只有"$38\,^{\circ}\mathrm{C}$——相对于新设定值，你"偏低"了，于是身体执行升温程序：寒战产热、血管收缩、你主观想盖被子。等退烧时设定值调回 $37\,^{\circ}\mathrm{C}$，实际体温又"偏高"了，于是你大汗淋漓地散热。忽冷忽热的发烧体验，就是设定值先动、体温后追的两段纠偏过程。

那谁在动设定值？是免疫细胞释放的信号分子（细胞因子，第 61 章的主角们），它们随血液到达下丘脑，命令局部细胞合成前列腺素，把恒温器的刻度往上拨。演化为什么保留这套程序？多数证据指向：适度升温能抑制一部分病原体的繁殖速度，同时让免疫细胞跑得更勤快。所以发烧更像一次"主动提高车间温度赶工灭菌"的操作，而不是事故。

当然，这不等于发烧一律不用管。温度太高（比如超过 $40\,^{\circ}\mathrm{C}$）、持续太久，或发生在婴幼儿、老人、有基础病的人身上，风险会压过收益——退烧药的作用正是把设定值拨回去。是否用药、何时就医，取决于具体的人和具体的情况，这是医生要权衡的事，不是一个固定数字能回答的。记住机制，比记住"几度该吃药"有用得多。

\section{模型在哪里会失灵}

稳态这张图纸好用到什么程度？也好用到必须说清楚它的边界。

\begin{itemize}[itemsep=2pt]
\item \textbf{纠偏能力有上限}。高温高湿环境里，汗蒸发不掉，散热这条路被物理堵死，体温会被产热一路顶上去——中暑就是散热纠偏彻底输掉的状态，这时身体不是"不想调"，是调不动。长时间严寒、严重腹泻脱水同理。
\item \textbf{零件会坏}。胰岛 $\beta$ 细胞被破坏或身体对胰岛素不敏感，血糖这条回路就缺了控制中心，血糖长期居高不下——这就是糖尿病的本质：一个坏掉的负反馈回路。治疗思路也由此而来：缺胰岛素的补充胰岛素（补上执行指令），不敏感的去改善敏感性或用药帮忙降糖。
\item \textbf{设定值不是普世常数}。所谓"正常范围"来自大量人群的统计，个体之间有差异：有人基础体温常年 $36.0\,^{\circ}\mathrm{C}$，有人 $36.8\,^{\circ}\mathrm{C}$，都在各自的健康轨道上。年龄也在改范围——老年人的体温调节和口渴感都会变迟钝，所以他们更怕热浪；婴幼儿的回路还没调试成熟，体温波动更大。女性的体温还随月经周期有规律的起伏。
\item \textbf{"稳"的范围本身会被改写}。长期住在高原的人，血液里的红细胞数量会上调，去适应低氧——这不是把血氧拉回平原标准，而是把整条生产线搬了新家。有人用"应变稳态"（allostasis）来描述这种"按环境重设工作点"的长期策略。它提醒我们：稳态的"设定值"也是演化和适应的产物，不是刻在石头上的。
\end{itemize}

还有一层更诚实的边界：这一章把每个变量分开讲——体温一条回路、血糖一条回路、pH 一条回路。真实的身体里这些回路盘根错节：血糖崩溃会影响体温，脱水会牵连血压和肾功能，发烧时心跳和呼吸全都跟着变。分开讲是为了让你看清图纸，不是身体真的分间办公。

\begin{wuqu}
"出汗能排毒，多出汗等于给身体做大扫除。"——汗液 $99\,\%$ 以上是水和少量钠、钾、尿素，它存在的第一使命是散热：每蒸发一克汗带走约 $2.4$ 千焦热量，这是高效的物理降温。至于"毒"，真正的代谢废物和外来的药物、酒精，走的是肝转化、肾排泄这条流水线，汗液里那点儿尿素浓度跟尿液比完全不够看。刻意捂汗、高温下逼自己大量出汗，不仅排不了毒，还会丢掉水和盐——而水和盐恰恰是内环境拼命维稳的对象，丢了它们，你是给纠偏系统\textbf{制造}新麻烦。想"排毒"，好好喝水、让肝肾正常工作，比出一身汗管用得多。
\end{wuqu}

\section{动手试一试：看负反馈现场直播}

\begin{shiyan}
\textbf{材料}：一只带秒表的手机；一个安静的座位。

\textbf{步骤}：
\begin{enumerate}[itemsep=1pt]
\item 安静坐 $5$ 分钟，然后用手指按住手腕内侧（拇指那侧）的桡动脉，数 $30$ 秒脉搏，乘 $2$，记下安静心率。
\item 原地做 $1$ 分钟中强度运动（深蹲或开合跳），停下后立刻测 $30$ 秒脉搏，记下运动后心率。
\item 之后每隔 $1$ 分钟测一次，连续测 $4$ 到 $5$ 次，把数据连成一条折线。
\end{enumerate}

\textbf{观察点}：运动后心率冲到了多少？几分钟后回到安静值附近？曲线是"断崖式"掉下来，还是一点点逼近？想一想：心率升高是谁下的指令、为了满足什么偏差（肌肉耗氧猛增）？恢复过程对应负反馈的哪个环节？如果恢复得很慢，可能有哪些原因？

\textbf{安全提示}：运动强度以"微喘但能说话"为限；头晕、胸闷立刻停下休息。有心脏疾病或医生叮嘱不宜运动的同学，只做安静心率的测量和曲线分析即可，不要勉强运动。
\end{shiyan}

\section{回到那支体温计}

现在再看体温计上那个雷打不动的 $36$ 点几度，你读到的已经不是"身体很稳定"，而是一份纠偏系统的工作简报：在你夹体温计的五分钟里，你的汗腺在微调开度，你的皮肤血管在微调口径，你的肝脏在不动声色地拆糖原放糖，你的肾脏在逐滴决定回收多少水，你的下丘脑像调度塔台一样同时盯着十几路信号。体温计上的"不动"，是无数分子在动。

这也可以回答一个更隐蔽的问题：为什么体检报告上每项指标后面都印着一个"参考范围"，而不是一个标准值？因为医学检测的本质就是\textbf{抽查内环境}——血糖、血钠、尿酸、转氨酶，全是在问"你的纠偏系统还 hold 得住吗"。单项略出范围不必立刻惊慌（范围来自人群统计，也有个体波动和测量误差），但持续、大幅地偏离，往往意味着某条回路的传感器、控制中心或执行器出了故障，值得医生顺着回路去查。下次看体检报告，试着把每项指标对应到一条纠偏回路，你读的就不再是天书。

\section{那支纠偏队伍里，还有一支特殊兵种}

到这里，这一章讲的都是"对数字的纠偏"：温度、浓度、酸碱度，全是能写进化验单的连续变量。但身体面对的威胁里，有一类不是数字，而是\textbf{活的入侵者}——细菌、病毒、寄生虫，它们不按负反馈出牌，它们会自我复制、会变异、会跟你打游击。对付它们，"拉回设定值"的思路完全不够用，你需要的是一支能\textbf{识别敌我、定点清除、记住仇家}的部队。

这就是下一章的主角。发烧时调高设定值的那些信号分子，正是这支部队的传令兵；疫苗为什么能防病，答案也藏在这支部队的"记性"里。稳态管住了内环境的"物理化学"，而免疫管住了内环境的"治安"——两座系统在下丘脑和血液里随时互通消息。第 61 章，我们去见见这支会学习的军队。

\begin{tiaozhan}
负反馈可以用一行代数看穿。设某变量偏离设定值的误差为 $E$，纠偏系统每一轮能消除误差的 $k$ 倍（$0<k<1$），同时环境每一轮新添一个大小为 $D$ 的扰动。那么下一轮的误差是 $E' = (1-k)E + D$。稳态时误差不再变化，即 $E = (1-k)E + D$，解得 $E = D/k$。这个简单式子说出三条深刻的话：第一，只要扰动 $D$ 持续存在，误差\textbf{永远回不到零}——稳态是"小误差下的动态平衡"，不是完美归零，这正是"不停纠偏"的数学版；第二，纠偏效率 $k$ 越高，残余误差越小，血糖、体温调得越"平"；第三，$k$ 接近 $0$（控制中心失灵）时误差爆炸——糖尿病就是 $k$ 塌方后的 $D/k$ 发散。真实的生理系统比这复杂：纠偏强度往往随误差大小非线性变化，还有延迟。延迟配上过强的纠偏会产生来回振荡——这解释了为什么有些病人的血糖、体温会像秋千一样荡。跳过这一段不影响阅读主线，但下次再听到"动态平衡"四个字，你脑子里可以浮现出这个 $E=D/k$。
\end{tiaozhan}

\section*{练一练}

\begin{sikaoti}
\item 画一张"负反馈三角色"流程图（传感器→控制中心→执行器→回到传感器），用"喝了一大瓶水之后，身体把血钠浓度拉回正常范围"为例，标出每个角色在这个例子里分别是谁、用什么信号通信。
\item 发烧初期人会打寒战、觉得冷。用"设定值先上调、体温后追赶"的思路，画出一条时间轴上的两条曲线（设定值曲线和实际体温曲线），标出发烧期、平台期、退烧期，并解释退烧时为什么大汗淋漓。
\item 剧烈运动时肌肉大量产热、产生 \chem{CO_2} 和乳酸。请列出身体同时启动的至少三条纠偏回路（提示：分别针对热量、\chem{CO_2}、\chem{H^{+}}），并指出每条回路的执行器和大致时间尺度。
\item 这一章说"稳态是 $E=D/k$ 附近的小幅震荡，不是纹丝不动"。据此反驳这个说法："健康人的血糖应该是一条直线，波动说明身体出了问题。"波动大到什么程度才真正值得担心？你的判断依据是什么？
\item 正反馈若不及时终止就是灾难。请以"分娩"和"血液凝固"为例，各找出大自然给正反馈循环安装的"终止开关"是什么；再想想：如果凝血正反馈没有刹车，血管里会发生什么？
\item 查一查你最近一次体检（或家人体检）报告上的三个指标及其参考范围，各写一句话说明：这个指标反映哪条纠偏回路的状态，传感器、控制中心、执行器分别可能是谁。
\end{sikaoti}
